% Compile from this directory: latexmk -pdf report.tex

\documentclass[12pt,a4paper]{article}
\usepackage{cmap}
\usepackage[utf8]{inputenc}
\usepackage[T2A]{fontenc}
\usepackage[english,russian]{babel}
\usepackage{amsmath,amssymb,bm}
\usepackage{graphicx,float,booktabs,geometry,array}
\usepackage[justification=raggedright,singlelinecheck=false]{caption}
\usepackage[expansion=false]{microtype}
\usepackage{xcolor}
\usepackage{hyperref}
\geometry{margin=2.2cm}
\setlength{\emergencystretch}{2em}
\hypersetup{colorlinks=true,linkcolor=black,citecolor=black,urlcolor=blue}
\graphicspath{{../results/report_figures/}}
\makeatletter
\setlength{\@fptop}{0pt}
\setlength{\@fpsep}{12pt}
\setlength{\@fpbot}{0pt plus 1fil}
\makeatother
\newcommand{\reportfigure}[3]{%
\begin{figure}[H]
\centering
\includegraphics[width=0.96\textwidth]{#1}
\caption{#2}
\label{#3}
\end{figure}
}
\title{Динамика и пространственная структура экситон-поляритонных конденсатов при нерезонансной накачке}
\author{Коростелев Никита}
\date{}
\begin{document}
\maketitle
\clearpage

\section{Цель работы}\label{sec:1}

В предыдущем отчёте рассматривались дисперсия и моды поляритонной микрополости. Здесь исследуется динамика конденсата при нерезонансной накачке: его возникновение, релаксация и пространственная структура.

Основные вопросы --- как обмен частицами с резервуаром определяет порог и время релаксации, когда допустимо исключить резервуар из динамического описания и как форма накачки влияет на плотность, токи и циркуляцию. Сначала исследуется однородная система и её сокращённые модели, затем --- двумерная динамика при гауссовой и кольцевой накачках.

Используются модель и параметры Alyatkin et al. $(2024)$ \cite{ref3}, но с более простыми профилями накачки. Параметры для микрополости из предыдущего отчёта не пересчитываются; экспериментальная геометрия статьи \cite{ref3} не воспроизводится.

\clearpage

\section{Физическая модель и параметры}\label{sec:2}

\subsection{Что описывают поле и резервуар}\label{sec:2.1}

Динамические переменные модели:

\begin{itemize}\raggedright\setlength{\itemsep}{3pt}\setlength{\parsep}{0pt}

\item $\psi(r,t)$ --- комплексное поле конденсата с нормировкой $n_{C}=|\psi|^{2}$. Его пространственная фаза определяет ток, а временная --- частоту поля.

\item $n_{R}(r,t)$ --- плотность активного некогерентного резервуара, который пополняет конденсат через стимулированное рассеяние. Резервуар описывается населённостью, без собственной когерентной фазы.

\end{itemize}

Здесь $r=(x,y)$ --- координата в плоскости микрополости. Мы используем одну компоненту поля $\psi$ и одну плотность резервуара, не учитывая отдельно спиновые компоненты поляритонов.

Далее стационарность плотности и стационарность комплексного поля рассматриваются отдельно: установившийся режим допускает $\psi(r,t)=\psi_{0}(r)\exp(-i \Omega t)$.

\subsection{Уравнения модели: смысл и происхождение слагаемых}\label{sec:2.2}

Используем открытую диссипативную модель Гросса--Питаевского (\foreignlanguage{english}{Wouters and Carusotto, 2007}) \cite{ref1} в параметризации Alyatkin et al. $(2024)$ \cite{ref3}:

\begin{equation}\label{eq:1}
\begin{aligned}
i\hbar\frac{\partial\psi}{\partial t}
&=\left[-\frac{\hbar^2}{2m}\nabla_\perp^2
+g_C|\psi|^2+g_R(n_R+\eta P)\right.\\
&\qquad\left.+\frac{i\hbar}{2}(Rn_R-\gamma_C)\right]\psi ,
\end{aligned}\end{equation}

\begin{equation}\label{eq:2}
\frac{\partial n_R}{\partial t}=P(\mathbf r)-\gamma_R n_R-Rn_R|\psi|^2.\end{equation}

Первое уравнение задаёт эволюцию поля конденсата, второе --- изменение количества частиц в резервуаре. Накачка $P(r)$ в рассматриваемых расчётах не меняется со временем. Оператор $\nabla_{\perp}^{2}=\partial_{x}^{2}+\partial_{y}^{2}$ учитывает изменение поля по двум координатам.

Члены уравнения \eqref{eq:1} удобно разобрать по отдельности:

\begin{itemize}\raggedright\setlength{\itemsep}{3pt}\setlength{\parsep}{0pt}

\item $-(\hbar^{2}/2m) \nabla_{\perp}^{2}$ --- кинетический оператор, соответствующий параболической дисперсии $E(k)=\hbar^{2} k^{2}/(2m)$ около минимума нижней поляритонной ветви. Замена $k\to -i \nabla$ даёт лапласиан и знак минус. Для плоской волны оператор возвращает положительную энергию $\hbar^{2} k^{2}/(2m)$.

\item $g_{C}|\psi|^{2}=g_{C} n_{C}$ --- сдвиг энергии за счёт взаимодействия поляритонов в приближении среднего поля. Локальное парное взаимодействие даёт плотность энергии $(g_{C}/2)n_{C}^{2}$; её производная по $n_{C}$ равна $g_{C} n_{C}$. Отсюда кубический член $g_{C}|\psi|^{2} \psi$ в уравнении поля. При $g_{C}>0$ взаимодействие отталкивающее.

\item $g_{R} n_{R}$ --- сдвиг энергии конденсата из-за взаимодействия с активным резервуаром. В принятом локальном приближении энергия взаимодействия на единицу площади равна $g_{R} n_{C} n_{R}$, поэтому энергия на добавленную частицу конденсата сдвигается на $g_{R} n_{R}$. Положительный $g_{R}$ создаёт отталкивающий профиль в областях высокой плотности резервуара.

\item $g_{R} \eta P$ --- дополнительный синий сдвиг от фоновых тёмных экситонов в параметризации Alyatkin et al. $(2024)$ \cite{ref3}. Их вклад представлен эффективной плотностью $n_{\rm bg}=\eta P$, пропорциональной заданной накачке; поэтому $\eta$ имеет размерность времени. Это феноменологическая параметризация фона, без отдельного динамического уравнения и без его вклада в усиление $R n_{R}$.

\item $+(i \hbar/2)R n_{R}$ --- стимулированное усиление. В уравнении населённости скорость поступления пропорциональна $R n_{R} n_{C}$: она зависит и от доступных частиц резервуара, и от заселённости конденсата. Поскольку $n_{C}=|\psi|^{2}$, соответствующая скорость роста амплитуды равна $R n_{R}/2$. Умножение уравнения для амплитуды на $i \hbar$ даёт указанный член в \eqref{eq:1}.

\item $-(i \hbar/2)\gamma_{C}$ --- потери при постоянной скорости распада населённости. Закон $\mathrm d n_{C}/\mathrm dt=-\gamma_{C} n_{C}$ соответствует амплитуде $\psi\propto \exp(-\gamma_{C} t/2)$, откуда множитель $1/2$ и отрицательный знак мнимого энергетического сдвига.

\end{itemize}

Различие между энергетическими сдвигами и усилением видно после деления \eqref{eq:1} на $i \hbar$. Вещественный локальный сдвиг $V$ даёт $\mathrm d \psi/\mathrm dt=-i V \psi/\hbar$ и вращает фазу; мнимый член $i \hbar G/2$ даёт $\mathrm d \psi/\mathrm dt=G \psi/2$ и меняет амплитуду. Здесь $G=R n_{R}-\gamma_{C}$. Кинетический оператор, помимо эволюции фазы, обеспечивает пространственное перераспределение плотности.

В уравнении резервуара \eqref{eq:2} используются следующие законы обмена:

\begin{itemize}\raggedright\setlength{\itemsep}{3pt}\setlength{\parsep}{0pt}

\item $+P(r)$ --- заданная скорость заполнения резервуара. Накачка нерезонансная, поэтому этот источник входит в уравнение населённости, а не как когерентная добавка к $\psi$.

\item $-\gamma_{R} n_{R}$ --- линейные потери резервуара: предполагается постоянная вероятность ухода частицы в единицу времени.

\item $-R n_{R}|\psi|^{2}=-R n_{R} n_{C}$ --- истощение резервуара тем же процессом, который даёт положительный член $+R n_{R} n_{C}$ в уравнении плотности конденсата. Такое совпадение коэффициентов обеспечивает баланс частиц при переходе между компонентами.

\end{itemize}

Таким образом, резервуар играет две роли: он питает конденсат частицами и одновременно создаёт отталкивающий энергетический сдвиг. Из-за этого максимум плотности конденсата не обязан совпадать с максимумом накачки.

\subsection{Уравнение непрерывности}\label{sec:2.3}

Чтобы перейти от поля к плотности, используем

\[
\frac{\partial n_C}{\partial t}=\psi^*\frac{\partial\psi}{\partial t}+\psi\frac{\partial\psi^*}{\partial t}.
\]

После подстановки \eqref{eq:1} вещественные локальные энергетические члены сокращаются, усиление и потери дают источник плотности, а кинетический член --- дивергенцию тока:

\begin{equation}\label{eq:3}
\frac{\partial n_C}{\partial t}+\nabla\cdot\mathbf j=(Rn_R-\gamma_C)n_C,\end{equation}

\begin{equation}\label{eq:4}
\mathbf j=\frac{\hbar}{m}\operatorname{Im}(\psi^*\nabla\psi).\end{equation}

Таким образом, плотность меняется как из-за локального обмена частицами, так и из-за их переноса. Даже при положительном усилении выходящий поток может уменьшать плотность в выбранной области. В однородной системе $\nabla\cdot(\mathbf j)=0$, и условие роста сводится к $R n_{R}>\gamma_{C}$.

\subsection{Единицы, константы и параметры}\label{sec:2.4}

Длины измеряются в $\mu\mathrm m$, время --- в $\mathrm{ps}$, энергия --- в $\mathrm{meV}$. Поскольку мы описываем двумерную плотность, $n_{C}$ и $n_{R}$ имеют размерность $\mu\mathrm m^{-2}$, поле $\psi$ --- $\mu\mathrm m^{-1}$, а $P$ --- $\mu\mathrm m^{-2}\,\mathrm{ps}^{-1}$.

Параметры модели:

\begin{center}\begin{tabular}{>{\raggedright\arraybackslash}p{0.17\textwidth}>{\raggedright\arraybackslash}p{0.27\textwidth}>{\raggedright\arraybackslash}p{0.44\textwidth}}\toprule

\textbf{Обозначение} & \textbf{Значение} & \textbf{Физический смысл}\\\midrule

$\hbar$ & $0.6582\,\mathrm{meV}\,\mathrm{ps}$ & Приведённая постоянная Планка\\

$m$ & $5.63 \times 10^{-5} m_{e}$ & Эффективная масса; $m_{e}$ --- масса электрона\\

$g_{C}$ & $0.0024\,\mathrm{meV}\,\mu\mathrm m^{2}$ & Взаимодействие внутри конденсата\\

$g_{R}$ & $0.0048\,\mathrm{meV}\,\mu\mathrm m^{2}$ & Энергетический сдвиг от резервуара\\

$R$ & $0.021\,\mu\mathrm m^{2}\,\mathrm{ps}^{-1}$ & Стимулированное рассеяние в конденсат\\

$\gamma_{C}$ & $0.182\,\mathrm{ps}^{-1}$ & Потери конденсата\\

$\gamma_{R}$ & $0.050\,\mathrm{ps}^{-1}$ & Потери резервуара\\

$\eta$ & $9.8\,\mathrm{ps}$ & Фоновая плотность $n_{\rm bg}=\eta P$; $\eta \gamma_{R}=0.49$\\

\bottomrule\end{tabular}\end{center}

Характерные времена потерь равны $1/\gamma_{C}=5.49\,\mathrm{ps}$ и $1/\gamma_{R}=20\,\mathrm{ps}$. Это времена затухания соответствующих населённостей без пополнения. Они не означают, что конденсат обязательно достигает стационарного состояния за $5\text{--}20\,\mathrm{ps}$: вблизи порога релаксация может быть гораздо медленнее.

Для кинетической энергии удобно ввести коэффициент
\[
K=\hbar^{2}/(2m)\simeq 0.677\,\mathrm{meV}\,\mu\mathrm m^{2}.
\]
Тогда энергия плоской волны с волновым вектором $k$ равна $K k^{2}$. Эта запись понадобится при обсуждении импульсных распределений.

Используемые параметры основаны на феномено\-логической модели и оценках Alyatkin et al. $(2024)$ \cite{ref3}. Не все они определялись независимыми прямыми измерениями: в частности, $R$ подбирался в рамках модели, а $g_{R}=2g_{C}$ задаёт принятое соотношение взаимодействий. Это важно учитывать при сопоставлении расчётов с экспериментом.

\subsection{\texorpdfstring{Откуда возникает порог $P_{\rm th}$}{Откуда возникает порог Pth}}\label{sec:2.5}

Порог можно понять ещё до полного решения стационарной задачи. Рассмотрим пространственно однородную накачку $P$ и момент, когда конденсат пока практически отсутствует: $n_{C}$ приблизительно равно нулю.

\begin{enumerate}\raggedright\setlength{\itemsep}{3pt}\setlength{\parsep}{0pt}

\item Если конденсат ещё не забирает частицы, в установившемся резервуаре поступление равно потерям: $P=\gamma_{R} n_{R}$. Поэтому $n_{R}=P/\gamma_{R}$.

\item Граница между затуханием и ростом малой амплитуды определяется равенством усиления и потерь: $R n_{R}=\gamma_{C}$.

\item Подставляя $n_{R}=P/\gamma_{R}$ во второе условие, получаем $R P/\gamma_{R}=\gamma_{C}$, то есть

\end{enumerate}

\begin{equation}\label{eq:5}
P_{\rm th}=\frac{\gamma_C\gamma_R}{R}=0.433333\ldots\;\mu\mathrm m^{-2}\,\mathrm{ps}^{-1},\end{equation}

Выражение \eqref{eq:5} относится к однородной системе. При локализованной накачке порог дополнительно зависит от пространственной формы моды и её выхода из области усиления; соответствующая задача рассматривается в разделе \ref{sec:8}.

Для дальнейшего анализа удобно определить $s=P/P_{\rm th}$ --- накачку в единицах этого порога. Для локализованного профиля будем писать $P(r)=P_{0} P_{\rm th} f(r)$, где максимальное значение $f$ равно единице. Здесь $P_{0}$ --- безразмерная высота профиля, а не полная мощность накачки.

\clearpage

\section{Пространственно однородные стационарные состояния и устойчивость}\label{sec:3}

\subsection{Переход от комплексного поля к двум плотностям}\label{sec:3.1}

Сначала рассмотрим самый простой случай: накачка однородна, а все величины одинаковы в каждой точке плоскости. Тогда пространственные производные поля равны нулю, и движение частиц между различными областями не влияет на динамику. Остаётся обмен частицами между резервуаром и конденсатом.

Уравнение для плотности конденсата уже получено из \eqref{eq:1} в виде \eqref{eq:3}. При отсутствии пространственных градиентов $\nabla\cdot(\mathbf j)=0$, поэтому

\begin{equation}\label{eq:6}
\frac{\mathrm dn_C}{\mathrm dt}=(Rn_R-\gamma_C)n_C,\end{equation}

Уравнение резервуара при этом не меняет своей формы:

\begin{equation}\label{eq:7}
\frac{\mathrm dn_R}{\mathrm dt}=P-\gamma_R n_R-Rn_Rn_C.\end{equation}

Эти уравнения удобно читать вместе. Если в резервуаре много частиц, множитель $R n_{R}-\gamma_{C}$ становится положительным, и конденсат растёт. Но рост конденсата увеличивает последний отрицательный член уравнения \eqref{eq:7}, то есть ускоряет истощение резервуара. Поэтому усиление не может бесконечно оставаться постоянным: возникает обратная связь, ограничивающая рост плотности.

Эта пара уравнений описывает изменение населённостей. Фаза поля в неё не входит, поскольку локальные энергетические сдвиги не изменяют $n_{C}$ напрямую.

\subsection{Как найти стационарные состояния}\label{sec:3.2}

Положим правые части \eqref{eq:6}--\eqref{eq:7} равными нулю и найдём стационарные плотности.

Условие $(R n_{R}-\gamma_{C})n_{C}=0$ выделяет две ветви: $n_{C}=0$ и ненулевую ветвь с $R n_{R}=\gamma_{C}$.

Ветвь без конденсата. Полагаем $n_{C}=0$. Тогда из второго уравнения следует $P-\gamma_{R} n_{R}=0$, откуда

\begin{equation}\label{eq:8}
n_C^{(0)}=0,\qquad n_R^{(0)}=\frac{P}{\gamma_R}.\end{equation}

Это решение математически существует при любой $P$. Но при достаточно сильной накачке оно окажется неустойчивым: даже маленькая ненулевая затравка конденсата начнёт расти.

Ветвь с конденсатом. Теперь $n_{C}>0$, поэтому обязательно $R n_{R}=\gamma_{C}$. Это означает, что резервуар должен поддерживать именно такую плотность, при которой усиление конденсата точно компенсирует его потери:

\[
n_R^*=\frac{\gamma_C}{R}.
\]

Подставим это значение в стационарное уравнение резервуара:

\[
0=P-\gamma_R\frac{\gamma_C}{R}-\gamma_C n_C^*.
\]

Учитывая определение $P_{\rm th}$ из \eqref{eq:5}, получаем:

\begin{equation}\label{eq:9}
n_R^*=\frac{\gamma_C}{R},\qquad n_C^*=\frac{P-P_{\rm th}}{\gamma_C},\qquad P>P_{\rm th}.\end{equation}

Так становится понятен смысл порога. Если $P<P_{\rm th}$, формула дала бы отрицательную плотность конденсата, что физически невозможно. При $P>P_{\rm th}$ появляется допустимое решение с $n_{C}^{*}>0$.

Используя безразмерную накачку $s=P/P_{\rm th}$, получаем

\begin{equation}\label{eq:10}
n_C^*=\frac{\gamma_R}{R}(s-1),\qquad n_R^*=\frac{\gamma_C}{R},\end{equation}

Для выбранных параметров $\gamma_{R}/R=2.38095\,\mu\mathrm m^{-2}$, а $\gamma_{C}/R=8.66667\,\mu\mathrm m^{-2}$. Поэтому после порога плотность конденсата растёт линейно с $s$, а плотность активного резервуара остаётся постоянной. Это называют закреплением, или насыщением, резервуара: дополнительная накачка в стационарном состоянии прежде всего переходит в увеличение $n_{C}$.

Однако стационарность плотности не означает, что сама комплексная функция $\psi$ постоянна. При $n_{C}^{*}>0$ она имеет вид

\begin{equation}\label{eq:11}
\begin{aligned}
\psi(t)&=\sqrt{n_C^*}\exp(-i\omega_*t+i\varphi_0),\\
\hbar\omega_*&=g_Cn_C^*+g_R(n_R^*+\eta P).
\end{aligned}\end{equation}

Здесь $\varphi_{0}$ --- произвольная начальная фаза, а $\omega_{*}$ --- частота её вращения. Энергия $\hbar \omega_{*}$ отсчитывается от выбранной опорной энергии модели и не равна полной оптической энергии поляритонного резонанса.

\subsection{Зачем нужна матрица Якоби и как её получить}\label{sec:3.3}

Чтобы выяснить, возвращается ли система к стационарному состоянию после малого отклонения, линеаризуем уравнения для плотностей.

Обозначим правые части уравнений \eqref{eq:6} и \eqref{eq:7} через две функции:

\[
\begin{aligned}F_C(n_C,n_R)&=(Rn_R-\gamma_C)n_C,\\F_R(n_C,n_R)&=P-\gamma_R n_R-Rn_Rn_C.\end{aligned}
\]

Пусть $(n_{C}^{*},n_{R}^{*})$ --- выбранная стационарная точка, а $\delta n_{C}$ и $\delta n_{R}$ --- небольшие отклонения от неё:

\[
\begin{aligned}n_C(t)&=n_C^*+\delta n_C(t),\\n_R(t)&=n_R^*+\delta n_R(t).\end{aligned}
\]

Подставляем их в $F_{C}$ и $F_{R}$ и оставляем только члены первой степени по $\delta n_{C}$ и $\delta n_{R}$. Постоянные части исчезают, потому что в стационарной точке $F_{C}=F_{R}=0$; произведения двух малых отклонений отбрасываем. Получается система

\[
\frac{\mathrm d}{\mathrm dt}\begin{pmatrix}\delta n_C\\\delta n_R\end{pmatrix}=J\begin{pmatrix}\delta n_C\\\delta n_R\end{pmatrix}.
\]

Матрица $J$ состоит из частных производных правых частей по $n_{C}$ и $n_{R}$. Например, $\partial F_{C}/\partial n_{R}=R n_{C}$ показывает, насколько изменение резервуара меняет скорость роста конденсата. Все четыре производные дают

\begin{equation}\label{eq:12}
J(\mathbf u)=\begin{pmatrix}Rn_R-\gamma_C&Rn_C\\-Rn_R&-\gamma_R-Rn_C\end{pmatrix}.\end{equation}

где $u=(n_{C},n_{R})^{T}$; для анализа выбранного состояния $J$ вычисляется в его стационарной точке.

Для анализа устойчивости ищем отклонения вида $\delta n(t)=v \exp(\lambda t)$, где $v$ --- постоянный вектор. Подстановка даёт $J v=\lambda v$: число $\lambda$ должно быть собственным значением $J$. Если его действительная часть отрицательна, такое возмущение затухает; если положительна --- растёт. Нулевая действительная часть требует отдельного рассмотрения.

\subsection{Устойчивость состояния без конденсата}\label{sec:3.4}

Начнём с решения \eqref{eq:8}: $n_{C}=0$ и $n_{R}=P/\gamma_{R}$. Подставляя его в матрицу \eqref{eq:12}, получаем треугольную матрицу

\[
J_0=\begin{pmatrix}RP/\gamma_R-\gamma_C&0\\-RP/\gamma_R&-\gamma_R\end{pmatrix}.
\]

Для треугольной матрицы собственные значения равны её диагональным элементам. С учётом $s=P/P_{\rm th}$ и $P_{\rm th}=\gamma_{C} \gamma_{R}/R$ это

\begin{equation}\label{eq:13}
\lambda_C=\gamma_C(s-1),\qquad\lambda_R=-\gamma_R.\end{equation}

Значение $\lambda_{R}$ всегда отрицательно: если резервуар немного отклонить от его стационарной плотности при отсутствующем конденсате, это отклонение само по себе затухает.

Собственное значение $\lambda_{C}$ меняет знак при $s=1$:

\begin{itemize}\raggedright\setlength{\itemsep}{3pt}\setlength{\parsep}{0pt}

\item при $s<1$ малая плотность конденсата затухает, и ветвь без конденсата устойчива к рассмотренным однородным возмущениям;

\item при $s>1$ она растёт, и ветвь без конденсата становится неустойчивой;

\item при $s=1$ линейная скорость роста обращается в нуль --- это порог, найденный ранее из баланса усиления и потерь.

\end{itemize}

Именно эту смену знака показывает рис.~\ref{fig:1}(b). На нём изображено наибольшее собственное значение $\max(\lambda_{C},\lambda_{R})$. При слабой накачке максимумом может оказаться постоянное число $-\gamma_{R}$, поэтому на графике появляется горизонтальный участок. Это не дополнительный физический порог, а результат выбора большего из двух собственных значений.

Важно различать рост поля и рост плотности. Если амплитуда $\psi$ растёт как $\exp(\sigma t)$, то плотность $|\psi|^{2}$ растёт как $\exp(2 \sigma t)$. Поэтому показатель $\lambda_{C}$ для плотности вдвое больше показателя роста амплитуды соответствующей конденсатной моды. Это различие понадобится при обсуждении двумерных порогов в разделе \ref{sec:8}.

\reportfigure{01_homogeneous_threshold.png}{Порог конденсации в пространственно однородной модели и устойчивость ветви без конденсата. (a) Стационарные $n_{C}$ и $n_{R}$ в зависимости от $s=P/P_{\rm th}$; выше $s=1$ плотность $n_{C}$ растёт линейно, а $n_{R}$ остаётся закреплённой. (b) Наибольшее собственное значение матрицы Якоби для двух плотностей на ветви с нулевой плотностью конденсата. Штриховая линия обозначает аналитический порог.}{fig:1}

\subsection{Устойчивость ненулевой ветви и колебательная релаксация}\label{sec:3.5}

Для ветви с конденсатом выполняются $R n_{R}^{*}=\gamma_{C}$ и $R n_{C}^{*}=\gamma_{R}(s-1)$. Поэтому матрица \eqref{eq:12} в стационарной точке упрощается:

\[
J_*=\begin{pmatrix}0&\gamma_R(s-1)\\-\gamma_C&-\gamma_Rs\end{pmatrix}.
\]

Здесь особенно хорошо видна обратная связь: увеличение плотности резервуара усиливает конденсат (верхний правый элемент положителен), а увеличение плотности конденсата быстрее истощает резервуар (нижний левый элемент отрицателен).

Ищем собственные значения из условия $\det(J_{*}-\lambda I)=0$. Для матрицы $2 \times 2$ это квадратное уравнение

\[
\lambda^2+\gamma_Rs\lambda+\gamma_C\gamma_R(s-1)=0.
\]

Решая его, получаем

\begin{equation}\label{eq:14}
\lambda_\pm=-\frac{\gamma_Rs}{2}\pm\frac12\sqrt{\gamma_R^2s^2-4\gamma_C\gamma_R(s-1)}.\end{equation}

У обеих мод действительные части отрицательны при $s>1$. Это также следует из двух свойств матрицы: её след равен $-\gamma_{R} s<0$, а определитель --- $\gamma_{C} \gamma_{R}(s-1)>0$. Следовательно, ненулевая стационарная ветвь линейно устойчива к малым пространственно однородным возмущениям плотностей. Такой расчёт не доказывает устойчивость к произвольным двумерным возмущениям поля.

Тип возвращения к стационарному состоянию зависит от выражения под корнем в \eqref{eq:14}:

\begin{itemize}\raggedright\setlength{\itemsep}{3pt}\setlength{\parsep}{0pt}

\item если оно положительно, собственные значения вещественные, и малые отклонения затухают без колебаний;

\item если отрицательно, значения образуют комплексно-сопряжённую пару, и плотности совершают затухающие колебания.

\end{itemize}

Для выбранных параметров колебательная область приблизительно равна $1.0801<s<13.4799$. Например, при $s=1.5$

\begin{equation}\label{eq:15}
\lambda_\pm=-0.0375\pm0.0561\,i\;\mathrm{ps}^{-1}.\end{equation}

Число $0.0375\,\mathrm{ps}^{-1}$ задаёт скорость затухания малых отклонений, а $0.0561\,\mathrm{ps}^{-1}$ --- угловую частоту колебаний. Физически они возникают из-за запаздывающей обратной связи: рост конденсата истощает резервуар, усиление ослабевает, затем резервуар восстанавливается за счёт накачки. Граница между колебательным и неколебательным режимами определяется дискриминантом в \eqref{eq:14}.

Рисунок~\ref{fig:2}(a) показывает обе части собственных значений. На рис.~\ref{fig:2}(b) изображена уже численная траектория $(n_{C}(t),n_{R}(t))$ при $s=1.5$. Это фазовая траектория в пространстве двух плотностей, а не график оптической фазы $\psi$. Она иллюстрирует, как конденсат и резервуар вместе приближаются к стационарному состоянию.

При этом в уравнениях \eqref{eq:6}--\eqref{eq:7} отсутствует общая фаза комплексного поля. Её произвольное постоянное смещение не влияет на плотности и соответствует отдельной нейтральной фазовой моде; приведённые собственные значения описывают только плотностные возмущения.

\reportfigure{02_condensed_state_stability.png}{Устойчивость и релаксация однородной конденсатной ветви. (a) Действительные и мнимые части двух собственных значений системы «плотность конденсата --- резервуар» при $s>1$; комплексные значения соответствуют затухающим колебаниям населённостей. (b) Траектория $(n_{C},n_{R})$ при $s=1.5$ с указанием направления времени, начальной точки и стационарного состояния. Глобальная фазовая мода не рассматривается.}{fig:2}

\clearpage

\section{Пространственно однородная динамика и сокращённые модели}\label{sec:4}

Стационарные решения не определяют переходную динамику. Даже при одинаковых конечных плотностях модели могут давать разные времена роста, выбросы плотности и затухающие колебания. В этом разделе полная система сравнивается с двумя последовательными сокращениями: адиабатическим исключением резервуара и последующим кубическим разложением по полю.

Первое сокращение устраняет запаздывание резервуара; второе дополнительно приближает зависимость насыщения усиления. Их условия применимости и ошибки рассматриваются раздельно.

\subsection{Адиабатическое исключение резервуара}\label{sec:4.1}

Перепишем уравнение \eqref{eq:7}, собрав все потери резервуара в один множитель:

\[
\frac{\mathrm dn_R}{\mathrm dt}=P-(\gamma_R+Rn_C)n_R.
\]

Если на короткое время считать $n_{C}$ постоянной, резервуар стремится к значению, при котором правая часть равна нулю. Отсюда сразу получается

\begin{equation}\label{eq:16}
n_R^{(\rm ad)}(n_C)=\frac{P}{\gamma_R+Rn_C}.\end{equation}

В адиабатическом приближении динамическое уравнение резервуара заменяется зависимостью \eqref{eq:16}, вычисляемой при текущей $n_{C}$. Это предполагает, что отклонение $n_{R}$ от этой зависимости успевает затухнуть до заметного изменения конденсата.

В знаменателе складываются собственные потери $\gamma_{R}$ и истощение резервуара $R n_{C}$. Поэтому с ростом конденсата уменьшаются $n_{R}^{(\rm ad)}$ и усиление $R n_{R}$: насыщение следует из исходного баланса частиц.

Чтобы понять условие применимости, рассмотрим небольшое отклонение резервуара от \eqref{eq:16} при временно фиксированной $n_{C}$. Оно удовлетворяет уравнению $d(\delta n_{R})/\mathrm dt=-(\gamma_{R}+R n_{C})\delta n_{R}$ и затухает за характерное время

\[
\tau_R=\frac{1}{\gamma_R+Rn_C}.
\]

Адиабатическое исключение оправдано, если за это время плотность конденсата меняется мало. Если же конденсат и резервуар существенно меняются за сопоставимые времена, резервуар запаздывает и замена его динамики формулой \eqref{eq:16} может быть неточной. При начальном резервуаре, далёком от \eqref{eq:16}, нужно также учитывать начальный этап его подстройки.

При меняющейся $n_{C}$ зависимость $n_{R}^{(\rm ad)}(n_{C}(t))$ тоже меняется. Обнуление правой части уравнения резервуара задаёт её мгновенное равновесное значение, а не утверждает точное равенство $\mathrm d n_{R}^{(\rm ad)}/\mathrm dt=0$ вдоль всей траектории.

\subsection{Неразложенная адиабатическая модель: стационарность и динамика}\label{sec:4.2}

Разделим однородное уравнение поля \eqref{eq:1} на $i \hbar$ и вместо $n_{R}$ подставим выражение \eqref{eq:16}. Получится одно уравнение для комплексного поля:

\begin{equation}\label{eq:17}
\begin{aligned}
\frac{\mathrm d\psi}{\mathrm dt}
&=\left[\frac12\left(\frac{RP}{\gamma_R+Rn_C}-\gamma_C\right)\right.\\
&\qquad\left.-\frac{i}{\hbar}\left\{g_Cn_C+
g_R\left(\frac{P}{\gamma_R+Rn_C}+\eta P\right)\right\}\right]\psi,\\
n_C&=|\psi|^2 .
\end{aligned}\end{equation}

В первой круглой скобке стоит действительная скорость изменения амплитуды: усиление из резервуара минус потери. Выражение в фигурных скобках, умноженное на $-i/\hbar$, задаёт изменение фазы через энергию взаимодействий.

Зависимость $P/(\gamma_{R}+R n_{C})$ пока сохранена целиком. Поэтому будем называть \eqref{eq:17} неразложенной адиабатической моделью. Для неё не требуется малость $n_{C}$: отдельное ограничение здесь связано со скоростью подстройки резервуара.

В ненулевом стационарном состоянии амплитуда перестаёт расти. Действительная часть коэффициента в \eqref{eq:17} обращается в нуль:

\[
\frac{RP}{\gamma_R+Rn_C}=\gamma_C.
\]

Решение этого равенства даёт ту же плотность $n_{C}^{*}=(\gamma_{R}/R)(s-1)$, что и полная модель. Плотность резервуара также совпадает: $n_{R}^{*}=\gamma_{C}/R$. Состояние без конденсата воспроизводится аналогично.

Но совпадение конечных плотностей не означает совпадения пути к ним. В полной модели $n_{R}$ --- самостоятельная переменная с собственной начальной величиной и конечным временем отклика. В \eqref{eq:17} резервуар в каждый момент однозначно задан текущей $n_{C}$. Поэтому сокращённая модель может верно предсказать конечный уровень и одновременно неточно описать переходный выброс или затухающие колебания.

\subsection{Вывод уравнения Стюарта–Ландау}\label{sec:4.3}

Следующий шаг --- упростить дробь в \eqref{eq:16}. Для этого вынесем $\gamma_{R}$ из знаменателя и введём безразмерную величину $x=R n_{C}/\gamma_{R}$. Тогда $n_{R}=(P/\gamma_{R})/(1+x)$.

Если $x$ много меньше единицы, можно использовать разложение $1/(1+x)=1-x+O(x^{2})$:

\begin{equation}\label{eq:18}
n_R=\frac{P}{\gamma_R}[1-x+O(x^2)],\qquad x=\frac{Rn_C}{\gamma_R}\ll1.\end{equation}

Условие $x\ll 1$ означает, что истощение резервуара конденсатом мало по сравнению с собственными потерями резервуара. Оно ничего само по себе не говорит о том, успевает ли резервуар следовать за конденсатом. Быстрая подстройка и слабое истощение --- два отдельных требования.

Подставим \eqref{eq:18} в \eqref{eq:17}. В коэффициенте перед $\psi$ оставим постоянные члены и члены, пропорциональные $n_{C}=|\psi|^{2}$. Получаем

\begin{equation}\label{eq:19}
\frac{\mathrm d\psi}{\mathrm dt}=[\mu-\beta|\psi|^2-i(\omega_0+\nu|\psi|^2)]\psi,\end{equation}

Это уравнение Стюарта--Ландау для комплексной амплитуды с коэффициентами

\begin{equation}\label{eq:20}
\begin{aligned}
\mu&=\frac12\left(\frac{RP}{\gamma_R}-\gamma_C\right),&
\beta&=\frac{PR^2}{2\gamma_R^2},\\
\omega_0&=\frac{g_R(P/\gamma_R+\eta P)}{\hbar},&
\nu&=\frac{g_C-g_RPR/\gamma_R^2}{\hbar}.
\end{aligned}\end{equation}

Разложение до первой степени $n_{C}$ в коэффициенте перед $\psi$ оставляет члены до $|\psi|^{2} \psi$; следующий отброшенный порядок в уравнении поля --- $|\psi|^{4} \psi$.

Коэффициенты \eqref{eq:20} можно читать следующим образом:

\begin{itemize}\raggedright\setlength{\itemsep}{3pt}\setlength{\parsep}{0pt}

\item $\mu$ --- линейная скорость роста амплитуды слабого поля. Она положительна выше порога и отрицательна ниже него.

\item $\beta>0$ --- коэффициент насыщения усиления. Член $-\beta|\psi|^{2}$ уменьшает скорость роста по мере увеличения плотности.

\item $\omega_{0}$ --- частота вращения фазы в пределе слабого поля, обусловленная резервуарным и фоновым энергетическими сдвигами.

\item $\nu$ --- изменение этой частоты при увеличении плотности. В нём складываются собственное взаимодействие конденсата и уменьшение резервуарного сдвига из-за истощения резервуара.

\end{itemize}

В частности, $\hbar \nu$ содержит разность собственного взаимодействия конденсата и уменьшения резервуарного сдвига при истощении.

Положив $\psi=\sqrt{n_{C}}\exp(i \varphi)$, можно отдельно записать смысл действительной и мнимой частей \eqref{eq:19}:

\[
\begin{aligned}\frac{\mathrm dn_C}{\mathrm dt}&=2(\mu-\beta n_C)n_C,\\\frac{\mathrm d\varphi}{\mathrm dt}&=-(\omega_0+\nu n_C).\end{aligned}
\]

Таким образом, $\mu$ и $\beta$ задают рост и насыщение плотности, а $\omega_{0}$ и $\nu$ --- эволюцию фазы. Фоновый член $\eta P$ входит в $\omega_{0}$ и не меняет однородный порог усиления. Это утверждение относится к однородной модели; в пространственной задаче энергетический профиль влияет на форму мод и тем самым может влиять на порог.

Итак, для адиабатического исключения нужна быстрая подстройка резервуара, а для кубического разложения --- дополнительно $R n_{C}/\gamma_{R}\ll 1$. Эти условия независимы и не обеспечены автоматически для выбранных параметров. Поэтому обе редукции далее сравниваются с полной динамикой. Ограничения адиабатического описания обсуждаются Bobrovska и Matuszewski $(2015)$ \cite{ref2}.

\subsection{Стационарная ошибка кубического приближения}\label{sec:4.4}

Из уравнения для $n_{C}$ после \eqref{eq:20} видно, что ненулевая стационарная плотность определяется условием $\mu-\beta n_{C}=0$. Подставляя коэффициенты \eqref{eq:20} и выражая $P$ через $s=P/P_{\rm th}$, получаем

\begin{equation}\label{eq:21}
n_C^{(\rm SL,*)}=\frac{\mu}{\beta}=\frac{\gamma_R}{R}\frac{s-1}{s},\qquad s>1.\end{equation}

В полной модели стационарная плотность равна $n_{C}^{*}=(\gamma_{R}/R)(s-1)$. В кубическом приближении появляется дополнительный делитель $s$. Поэтому выше порога кубическая модель предсказывает меньшую плотность.

Относительное расхождение с полной стационарной ветвью равно

\begin{equation}\label{eq:22}
\varepsilon_{\rm ss}^{(\rm SL)}=\frac{|n_C^{(\rm SL,*)}-n_C^*|}{n_C^*}=1-\frac1s,\qquad s>1.\end{equation}

При $s=1.05$ ошибка составляет $4.8\%$, при $s=1.2$ --- $16.7\%$, при $s=1.5$ --- $33.3\%$. Эти числа относятся именно к коэффициентам \eqref{eq:20}, полученным нашим разложением. Они не являются универсальной ошибкой любого уравнения Стюарта--Ландау.

Связь с областью применимости особенно проста в стационарном состоянии полной модели: параметр разложения $x^{*}=R n_{C}^{*}/\gamma_{R}=s-1$. Вблизи порога он мал, а при удалении от порога малость теряется. Поэтому ухудшение кубического приближения закономерно.

У неразложенной адиабатической модели стационарное расхождение равно нулю аналитически. Это следствие одинаковых уравнений стационарного баланса, а не результат особенно точного численного интегрирования.

\subsection{Сравнение временных траекторий}\label{sec:4.5}

На рис.~\ref{fig:3} сопоставлены три модели при $s=0.8, 1.05, 1.2$ и $1.5$. Начальное поле везде одинаково: $\psi(0)=0.01\,\mu\mathrm m^{-1}$, то есть $n_{C}(0)=10^{-4}\,\mu\mathrm m^{-2}$.

В полной модели начальный резервуар выбран на адиабатической зависимости: $n_{R}(0)=P/[\gamma_{R}+R n_{C}(0)]$. Это устраняет начальное отклонение резервуара от \eqref{eq:16}, но не заставляет его оставаться на ней в дальнейшем. Когда $n_{C}$ начинает меняться, конечное время отклика снова становится существенным.

Все три однородные модели интегрируются адаптивным методом Рунге--Кутты восьмого порядка DOP853 (\texttt{scipy.integrate.solve\_ivp}). Относительный допуск равен $2\times10^{-11}$, абсолютный --- $2\times10^{-13}$. В сокращённых моделях резервуар не интегрируется отдельно: его кривые восстанавливаются из текущей плотности поля по соответствующим алгебраическим формулам. Реконструкция резервуара первого порядка для кубической модели ниже порога не показана.

Поведение кривых удобно разобрать по режимам:

\begin{itemize}\raggedright\setlength{\itemsep}{3pt}\setlength{\parsep}{0pt}

\item При $s=0.8$ усиления недостаточно, и конденсат затухает.

\item При $s=1.05$ превышение порога мало, поэтому начальный рост медленный.

\item При $s=1.2$ и $1.5$ становится заметным занижение конечной плотности кубической моделью, предсказанное формулой \eqref{eq:22}.

\item При $s=1.5$ полная система показывает переходный выброс плотности, сглаженный при мгновенном исключении резервуара.

\end{itemize}

Переходный выброс связан с запаздыванием резервуара: его усиление не успевает мгновенно следовать за меняющейся плотностью конденсата. Поэтому совпадение конечных уровней не гарантирует совпадения переходного процесса.

\reportfigure{03_matched_initial_dynamics.png}{Динамика однородной системы «конденсат--резервуар» при согласованных начальных условиях. Плотности конденсата (a--d) и резервуара (e--h) показаны для $s=0.8, 1.05, 1.2$ и $1.5$. Сравниваются полная модель, неразложенная адиабатическая редукция и кубическое разложение Стюарта--Ландау. Кривые резервуара в сокращённых моделях восстановлены алгебраически; кубическое приближение ниже порога не показано. Временные диапазоны панелей различаются.}{fig:3}

\subsection{Стационарная и динамическая ошибки}\label{sec:4.6}

На рис.~\ref{fig:4}(a) показана стационарная ошибка \eqref{eq:22}. Она отвечает на вопрос, совпадают ли конечные плотности. Для сравнения всей временной траектории используется другая величина:

\begin{equation}\label{eq:23}
\varepsilon_{\rm dyn}=\frac{1}{n_C^*}\sqrt{\frac1T\int_0^T[n_C^{(\rm approx)}(t)-n_C^{(\rm full)}(t)]^2\,\mathrm dt}.\end{equation}

Здесь $n_{C}^{(\rm approx)}$ --- плотность в одной из сокращённых моделей, $n_{C}^{(\rm full)}$ --- в полной, $T=1600\,\mathrm{ps}$ --- общий интервал сравнения, а $n_{C}^{*}$ --- стационарная плотность полной модели.

На рис.~\ref{fig:4}(b) ошибка рассчитана для двенадцати значений $s$ от $1.05$ до $1.5$. Она включает переходный и поздний участки, поэтому зависит от $T$: длительное совпадение конечных уровней может уменьшить среднюю ошибку, не улучшая описание ранней динамики. При $T\to\infty$ ошибка адиабатической модели стремится к нулю, а кубической --- к $1-1/s$. Формула \eqref{eq:23} используется только выше порога, где $n_{C}^{*}>0$.

\reportfigure{04_reduced_model_accuracy.png}{Точность сокращённых моделей однородного конденсата. (a) Аналитические относительные ошибки стационарной плотности: нуль для неразложенного адиабатического приближения и $1-1/s$ для заданного кубического разложения. (b) Усреднённое по времени среднеквадратичное отклонение плотности от полной модели на интервале $T=1600\,\mathrm{ps}$, рассчитанное для двенадцати значений накачки при согласованных начальных условиях.}{fig:4}

\clearpage

\section{Пространственная модель и численный метод}\label{sec:5}

\subsection{Геометрия накачки}\label{sec:5.1}

В двумерном расчёте решаются полные уравнения \eqref{eq:1}--\eqref{eq:2}: резервуар остаётся самостоятельной переменной. Рассматриваются центральное гауссово пятно и кольцо:

\begin{equation}\label{eq:24}
\begin{aligned}
P_{\rm spot}(r)&=P_0P_{\rm th}\exp(-r^2/d^2),\\
P_{\rm ring}(r)&=P_0P_{\rm th}\exp[-(r-r_0)^2/d^2],\\
r&=\sqrt{x^2+y^2},\qquad d=2.5\;\mu\mathrm m,\qquad r_0=10\;\mu\mathrm m.
\end{aligned}\end{equation}

Здесь $P_{0}$ --- максимальная накачка в единицах однородного порога, $d$ --- ширина профиля на уровне $1/e, r_{0}$ --- радиус кольца. Полная ширина на половине максимума равна $2d \sqrt{\ln 2}=4.16\,\mu\mathrm m$; для кольца это радиальная толщина профиля. $P_{0}$ характеризует высоту, а не интегральную мощность: при одинаковом пике кольцо накачивает большую площадь.

Для рис.~\ref{fig:5}--\ref{fig:7} выбраны $P_{0}=13.2$ у пятна и $P_{0}=3.16875$ у кольца --- примерно на $20\%$ выше соответствующих пространственных порогов из раздела \ref{sec:8}. Оба профиля радиально симметричны, но случайное начальное поле может нарушать эту симметрию.

\subsection{Расчёт и начальные условия}\label{sec:5.2}

\begin{center}\begin{tabular}{>{\raggedright\arraybackslash}p{0.4\textwidth}>{\raggedright\arraybackslash}p{0.23\textwidth}>{\raggedright\arraybackslash}p{0.23\textwidth}}\toprule

\textbf{Параметр} & \textbf{Гауссово пятно} & \textbf{Кольцо}\\\midrule

Область & $256 \times 256\,\mu\mathrm m^{2}$ & $256 \times 256\,\mu\mathrm m^{2}$\\

Сетка & $512 \times 512$ & $512 \times 512$\\

Пространственный шаг & $0.5\,\mu\mathrm m$ & $0.5\,\mu\mathrm m$\\

Временной шаг & $0.025\,\mathrm{ps}$ & $0.05\,\mathrm{ps}$\\

Конечное время & $500\,\mathrm{ps}$ & $2000\,\mathrm{ps}$\\

\bottomrule\end{tabular}\end{center}

Используется расщепление Стрэнга: кинетическая эволюция вычисляется в пространстве Фурье, а локальная эволюция поля и резервуара --- совместно методом Рунге--Кутты четвёртого порядка. Множитель кинетического полушага имеет вид

\begin{equation}\label{eq:25}
\exp\left[-i\frac{\hbar^2}{2m}(k_x^2+k_y^2)\frac{\Delta t}{2\hbar}\right].\end{equation}

Полная симметричная схема имеет второй порядок по времени. Граничные условия периодические, без поглощающего слоя; влияние границ контролируется отдельно. Подробности реализации собраны в приложении А.

Начальный резервуар пуст, а поле задаётся малой случайной затравкой, локализованной в области накачки. Она необходима, поскольку в детерминистском уравнении строго нулевое поле не начинает расти само. После задания начальных условий случайный источник в уравнения не добавляется. При сравнении сеток используется одна и та же непрерывная затравка.

\subsection{Контроль установления и точности}\label{sec:5.3}

Интегральные населённости определяются по всей расчётной области:

\begin{equation}\label{eq:26}
N_C(t)=\int n_C(\mathbf r,t)\,\mathrm d^2r,\qquad N_R(t)=\int n_R(\mathbf r,t)\,\mathrm d^2r.\end{equation}

Из \eqref{eq:2}--\eqref{eq:3} при периодических границах следует общий баланс

\begin{equation}\label{eq:27}
\frac{\mathrm d(N_C+N_R)}{\mathrm dt}=\int P(\mathbf r)\,\mathrm d^2r-\gamma_CN_C-\gamma_RN_R.\end{equation}

Стимулированное рассеяние переносит частицы между компонентами и сокращается в сумме. Поэтому проверяется выполнение \eqref{eq:27}, а не сохранение $N_{C}+N_{R}$ в открытой системе.

Установление режима оценивается на окнах $100\,\mathrm{ps}$ по населённостям и пространственной структуре. Основные допуски --- $0.5\%$ для изменений населённостей и $1\%$ для полей и плотностей. Это позволяет отличить неподвижное распределение от движущейся структуры с постоянным числом частиц. Точные критерии приведены в приложении А.

Чувствительность к дискретизации проверяется отдельно, сравнением независимых расчётов. Для пятна подтверждён второй порядок по времени, а основные расхождения населённостей укладываются в $0.5\%$. Проверка меньшей области дала предупреждение о влиянии границ; результаты нельзя распространять на произвольный размер области.

Для конечного кольцевого состояния при $2000\,\mathrm{ps}$ выполнены два запуска от $t=0$ с той же затравкой: на сетке $256^{2}$ вместо $512^{2}$ и с шагом $0.025$ вместо $0.05\,\mathrm{ps}$. Расхождения населённостей и профиля плотности не превышают $0.006\%$; воспроизводятся установившийся режим, суммарный заряд \mbox{$q=-3$} и угловой момент около $-3 \hbar$ на частицу. Их определения даны в разделе \ref{sec:6}, численные подробности --- в приложении А. Сравнение двух пространственных сеток подтверждает согласованность результатов, но не определяет пространственный порядок сходимости.

\clearpage

\section{Пространственная структура конденсата и потоки частиц}\label{sec:6}

\subsection{Плотность и энергетический профиль}\label{sec:6.1}

Резервуар одновременно усиливает поле и создаёт отталкивание. С учётом собственного взаимодействия конденсата локальный энергетический сдвиг равен

\begin{equation}\label{eq:28}
V_{\rm int}(\mathbf r,t)=g_Cn_C(\mathbf r,t)+g_R[n_R(\mathbf r,t)+\eta P(\mathbf r)].\end{equation}

Это действительная локальная часть оператора \eqref{eq:1}: первые два вклада зависят от решения, а фон $g_{R} \eta P$ задаётся накачкой. Установившаяся плотность определяется совместным действием этого профиля, усиления и кинетического переноса.

На рис.~\ref{fig:5} сопоставлены накачка, резервуар, конденсат и $V_{\rm int}$. При гауссовой накачке конденсат распространяется дальше узкой области резервуара: частицы выходят из области усиления и продолжают испытывать потери. При кольцевой накачке резервуарный потенциал и плотность конденсата имеют кольцевую структуру с малой плотностью в центре. Чтобы установить, сопровождается ли она циркуляцией, необходимо рассмотреть фазу и ток.

\reportfigure{05_spatial_profiles.png}{Скорость накачки $P$, плотность резервуара $n_{R}$, плотность конденсата $n_{C}$ и энергетический сдвиг среднего поля $V_{\rm int}=g_{C} n_{C}+g_{R}(n_{R}+\eta P)$. Верхний ряд: гауссово пятно, $P_{0}=13.2, t=500\,\mathrm{ps}$. Нижний ряд: кольцо, $P_{0}=3.16875, t=2000\,\mathrm{ps}$. Оба режима примерно на $20\%$ выше соответствующих пространственных порогов; $d=2.5\,\mu\mathrm m$ и $r_{0}=10\,\mu\mathrm m$. Для каждой колонки используется единая абсолютная цветовая шкала.}{fig:5}

\subsection{Фаза, ток и циркуляция}\label{sec:6.2}

Подстановка $\psi=\sqrt{n_{C}}\exp(i \varphi)$ в \eqref{eq:4} даёт

\begin{equation}\label{eq:29}
\mathbf j=\frac{\hbar}{m}n_C\nabla\varphi.\end{equation}

Градиент фазы задаёт скорость $v=\mathbf j/n_{C}=(\hbar/m)\nabla \varphi$, а плотность --- вклад области в ток. В нулях поля фаза не определена, поэтому там используется исходная формула \eqref{eq:4}.

Неподвижная плотность допускает ненулевой ток: из \eqref{eq:3} следует $\nabla\cdot(\mathbf j)=(R n_{R}-\gamma_{C})n_{C}$. Перенос может компенсировать локальные усиление и потери. На рис.~\ref{fig:6} у пятна ток преимущественно направлен наружу, а фазовые фронты почти концентрические. У кольца наблюдается закручивание фазы вокруг центральной области.

Для ненулевого поля на замкнутом контуре однозначность $\psi$ требует изменения фазы на целое число оборотов. Топологический заряд определяется как

\begin{equation}\label{eq:30}
q=\frac{1}{2\pi}\oint\nabla\varphi\cdot\mathrm d\mathbf l,\end{equation}

При положительном обходе против часовой стрелки \mbox{$q=-3$} означает изменение фазы на $-6 \pi$. Заряд вычисляется по накопленным фазовым приращениям, а не по числу цветовых полос на карте. Контур не должен пересекать нули поля или области с ненадёжно малой плотностью.

Для конечного поля при $2000\,\mathrm{ps}$ несколько контуров дают \mbox{$q=-3$}. Это суммарный заряд внутри них: он не различает одно ядро с зарядом $-3$ и несколько близких сингулярностей с той же суммой.

\subsection{Угловой момент}\label{sec:6.3}

Дополнительная характеристика циркуляции --- орбитальный угловой момент всего поля:

\begin{equation}\label{eq:31}
L_z=\int\psi^*[-i\hbar(x\partial_y-y\partial_x)]\psi\,\mathrm d^2r.\end{equation}

Поскольку $x \partial_{y}-y \partial_{x}=\partial_{\theta}$, для чистой угловой моды $\psi=F(r)\exp(i q \theta)$ получаем $L_{z}=q \hbar N_{C}$. Для произвольного поля нормированный момент $L_{z}/(\hbar N_{C})$ не обязан быть целым или совпадать с зарядом контура.

В найденном кольцевом состоянии он близок к $-3$. Вместе с фазовой циркуляцией это подтверждает выраженный вращательный поток при почти круговой плотности.

\reportfigure{06_particle_flow_and_phase.png}{Ток конденсата и распределение фазы. Карты плотности со стрелками тока (a,c) и карты фазы (b,d) соответствуют гауссовой и кольцевой накачкам при $500$ и $2000\,\mathrm{ps}$ соответственно. Штриховые окружности обозначают контуры диагностики заряда: конечное кольцевое состояние имеет суммарный заряд \mbox{$q=-3$} и средний орбитальный угловой момент около $-3 \hbar$ на частицу. Фаза в областях малой плотности скрыта.}{fig:6}

Многозарядные вихревые состояния обсуждаются Alperin и Berloff $(2021)$ \cite{ref4}. Наш результат относится к выбранной затравке: статистика выбора заряда и устойчивость к произвольным возмущениям не исследовались. В разделе \ref{sec:9} рассматривается, как затухает предшествующая этому состоянию угловая модуляция плотности.

\clearpage

\section{Распределения в импульсном пространстве}\label{sec:7}

\subsection{Определение спектра}\label{sec:7.1}

Импульсное распределение характеризует состав комплексного поля по плоским волнам. Преобразуется поле вместе с фазой, а затем берётся квадрат модуля:

\begin{equation}\label{eq:32}
\widetilde\psi(k_x,k_y,t)=\int\psi(x,y,t)\exp[-i(k_xx+k_yy)]\,\mathrm dx\,\mathrm dy.\end{equation}

\begin{equation}\label{eq:33}
I_k(k_x,k_y,t)=|\widetilde\psi(k_x,k_y,t)|^2.\end{equation}

Преобразование самой плотности дало бы спектр её пространственных неоднородностей, а не импульсный состав поля. Резервуар в \eqref{eq:32} не входит.

На рис.~\ref{fig:7} показано $\log_{10}[I_{k}/(1\,\mu\mathrm m^{2})]$ с общей абсолютной шкалой; панели не нормируются каждая на свой максимум. На осях отложен волновой вектор в $\mu\mathrm m^{-1}$, связанный с импульсом соотношением $p=\hbar k$. Дискретная нормировка приведена в приложении А.

\subsection{Интерпретация распределений}\label{sec:7.2}

Для гауссова пятна спектр образует кольцо при ненулевом $|k|$. Оно соответствует близким модулям волновых векторов при разных направлениях и согласуется с радиальным выходящим потоком. Это кольцо в пространстве импульсов, тогда как накачка в координатном пространстве сосредоточена в центральном пятне.

Панели $7$(a,b) относятся к областям $L=256$ и $192\,\mu\mathrm m$ при одинаковых пространственном и временном шагах. Сходство основных особенностей показывает их воспроизводимость при этом изменении области; подробный контроль границ обсуждается в приложении А.

Масштаб кинетической энергии задаётся дисперсией

\begin{equation}\label{eq:34}
E_{\rm kin}(k)=\frac{\hbar^2k^2}{2m}=Kk^2,\qquad K=0.677\;\mathrm{meV}\,\mu\mathrm m^2.\end{equation}

При $k\simeq 1\,\mu\mathrm m^{-1}$ энергия составляет около $0.677\,\mathrm{meV}$ --- величину порядка энергетического сдвига на рис.~\ref{fig:5}. Это оценка масштаба, а не локальное равенство $E_{\rm kin}=V_{\rm int}$.

Спектр кольцевого состояния более сложный и приблизительно концентрический: он зависит как от радиальной структуры, так и от угловой фазы. Однако $I_{k}$ не определяет знак циркуляции. У полей $F(r)\exp(\pm i q \theta)$ с одинаковой радиальной частью интенсивности совпадают. Поэтому вывод о \mbox{$q=-3$} основан на фазовых контурах, а спектр дополняет пространственную картину.

\reportfigure{07_momentum_distributions.png}{Интенсивность когерентного конденсата в импульсном пространстве $\log_{10}[I_{k}/(1\,\mu\mathrm m^{2})]$. (a,b) Гауссово пятно при $t=500\,\mathrm{ps}$ для $L=256$ и $192\,\mu\mathrm m$ соответственно, с одинаковыми пространственным и временным шагами. (c) Кольцо при $t=2000\,\mathrm{ps}$ и $L=256\,\mu\mathrm m$. Для всех панелей использована единая абсолютная логарифмическая шкала; пара гауссовых спектров позволяет сравнить размеры расчётной области, но не служит самостоятельной проверкой уточнения сетки.}{fig:7}

\clearpage

\section{Пороги конденсации при локализованной накачке}\label{sec:8}

\subsection{Линейная пространственная задача}\label{sec:8.1}

В однородной системе достаточно баланса $R n_{R}=\gamma_{C}$. При локализованной накачке часть моды находится за пределами усиления и испытывает потери. Поэтому положительный локальный рост около максимума $P$ ещё не гарантирует рост всей моды: необходимо учесть её пространственную форму.

На ветви без конденсата резервуар имеет стационарный профиль

\begin{equation}\label{eq:35}
\psi_0(\mathbf r)=0,\qquad n_{R,0}(\mathbf r)=\frac{P(\mathbf r)}{\gamma_R}.\end{equation}

Положим $\psi=\delta \psi$ и $n_{R}=n_{R,0}+\delta n_{R}$. В первом порядке сохраняются кинетика и действие стационарного резервуара. Собственное взаимодействие кубично по малому полю, произведение $\delta n_{R} \delta \psi$ имеет второй порядок, а истощение резервуара начинается с $|\delta \psi|^{2}$; эти члены в линейное уравнение поля не входят. После деления \eqref{eq:1} на $i \hbar$ получаем

\begin{equation}\label{eq:36}
\begin{aligned}
\partial_t\delta\psi
&=\frac{iK}{\hbar}\nabla_\perp^2\delta\psi
+\left[\frac12\left(\frac{RP(\mathbf r)}{\gamma_R}-\gamma_C\right)\right.\\
&\qquad\left.-\frac{ig_R}{\hbar}\left(\frac{P(\mathbf r)}{\gamma_R}+\eta P(\mathbf r)\right)\right]\delta\psi.
\end{aligned}\end{equation}

Возмущение резервуара отдельно затухает со скоростью $\gamma_{R}$. В \eqref{eq:36} энергетический профиль остаётся даже при нулевом конденсате: он влияет на форму моды и её перекрытие с усилением. Поэтому порог нельзя найти только по максимальному значению локальной скорости роста.

\subsection{Показатель роста и поиск порога}\label{sec:8.2}

Обозначив оператор \eqref{eq:36} через $A$, подставим $\delta \psi=u(r)\exp(\lambda t)$. Получаем $A u=\lambda u$, где $\sigma=\operatorname{Re} \lambda$ --- скорость роста амплитуды. Это пространственный аналог анализа якобиана: вместо двухкомпонентного вектора здесь ищется профиль поля.

При ненулевой начальной проекции на ведущую моду со временем выделяется максимальная $\operatorname{Re} \lambda$. Поскольку $N_{\rm lin}=\int|\delta \psi|^{2}d^{2}r$ меняется как $\exp(2 \sigma t)$,

\begin{equation}\label{eq:37}
\sigma=\frac12\frac{\mathrm d\ln N_{\rm lin}(t)}{\mathrm dt}.\end{equation}

Множитель $1/2$ переводит скорость изменения квадрата нормы в скорость роста амплитуды. При $\sigma<0$ поле затухает, при $\sigma>0$ растёт; переход через нуль задаёт порог.

Для каждого $P_{0}$ интегрируется \eqref{eq:36} с фиксированным резервуаром \eqref{eq:35}, а $\sigma$ оценивается после выделения ведущей моды. Значения накачки с разными знаками $\sigma$ ограничивают пороговый интервал. Это численные границы, а не статистический доверительный интервал. Проверки выделения моды и технический способ оценки роста описаны в приложении А.

\reportfigure{08_localized_condensation_thresholds.png}{Пороги линейной неустойчивости состояния без конденсата при локализованной накачке. Скорость роста амплитуды поля $\sigma$ рассчитывается при $n_{R,0}=P/\gamma_{R}$. (a,b) Диапазоны сканирования накачки для пятна и кольца; (c,d) увеличенные окрестности соответствующих порогов. Кружки обозначают $\sigma<0$ (устойчивость нулевого состояния), треугольники --- $\sigma>0$ (его неустойчивость). Серые полосы ограничивают $P_{0,c}$ интервалами $[10.9375,11.0000]$ для пятна и $[2.6250,2.640625]$ для кольца. Штриховые линии показывают однородный эталон $P_{0}=1$. Полосы представляют собой численные интервалы, а не доверительные интервалы.}{fig:8}

\subsection{Результаты}\label{sec:8.3}

Для пятна получен интервал $P_{0,c}\in [10.9375,11.0000]$, для кольца --- $[2.6250,2.640625]$. Оба порога выше однородного ориентира $P_{0}=1$: усиление должно компенсировать потери всей пространственной моды.

У кольца меньше пороговая высота профиля, но больше накачиваемая площадь. Поэтому эти числа не устанавливают меньшую суммарную пороговую мощность; для её сравнения нужен $\int P \mathrm d^{2}r$. Значения относятся к геометрии \eqref{eq:24} и выбранным параметрам.

Линейный анализ определяет начало роста слабого поля. Конечное насыщенное состояние и его циркуляция устанавливаются уже нелинейной динамикой.

\clearpage

\section{Переходная динамика конденсата при кольцевой накачке}\label{sec:9}

\subsection{От шести максимумов к почти круговой плотности}\label{sec:9.1}

При $1000\,\mathrm{ps}$ на кольце ещё видны шесть максимумов, перемещающихся по азимуту. Чтобы выяснить, сохраняется ли эта структура, та же траектория продолжена до $2000\,\mathrm{ps}$. На рис.~\ref{fig:9} населённости приближаются к $N_{C}\simeq 306$ и $N_{R}\simeq 6533$, а снимки плотности показывают, что движущиеся максимумы постепенно теряют контраст.

Изменение угловой структуры характеризуется гармоникой плотности

\begin{equation}\label{eq:38}
A_m(t)=\frac{2\left|\int_{\rm core} n_C(r,\theta,t)\exp(im\theta)\,\mathrm d^2r\right|}{\int_{\rm core} n_C(r,\theta,t)\,\mathrm d^2r}.\end{equation}

Здесь $\theta$ --- полярный угол, а $\mathrm{core}$ --- центральный квадрат $80 \times 80\,\mu\mathrm m^{2}$. Нормировка отделяет относительную модуляцию от изменения числа частиц. Множитель $2$ соответствует амплитуде косинусной модуляции при полном угловом интегрировании; особенности используемой области пояснены в приложении А.

Для шести максимумов выбирается $m=6$. Значения $A_{6}$ уменьшаются с $29.1\%$ при $1000\,\mathrm{ps}$ до $0.567\%$ при $1500\,\mathrm{ps}$ и $0.0107\%$ при $2000\,\mathrm{ps}$. Следовательно, шестилепестковая структура --- затухающий переходный процесс. Максимумы могут двигаться, пока падает их контраст; само наблюдение вращения ещё не означает существования постоянной вращающейся структуры.

Первое окно, удовлетворяющее критериям установления режима, --- $1500\text{--}1600\,\mathrm{ps}$. На позднем окне $1900\text{--}2000\,\mathrm{ps}$ изменение нормированной формы плотности составляет лишь $0.0234\%$. Это характеристика установления во времени; чувствительность к сетке и шагу проверена отдельно в разделе \ref{sec:5}.

\reportfigure{09_ring_transient_evolution.png}{Переходная эволюция при кольцевой накачке. (a,b) Интегральные населённости $N_{C}(t)$ и $N_{R}(t)$ от $0$ до $2000\,\mathrm{ps}$. (c--f) Плотность конденсата при $1000, 1200, 1500$ и $2000\,\mathrm{ps}$ с общей абсолютной цветовой шкалой. Шестикратная угловая модуляция плотности затухает, и возникает приблизительно осесимметричное стационарное распределение плотности; конечное состояние имеет суммарный заряд \mbox{$q=-3$} внутри рассматриваемых контуров (рис.~\ref{fig:6}).}{fig:9}

\subsection{Почему исчезновение лепестков не устраняет циркуляцию}\label{sec:9.2}

Угловая модуляция характеризует $|\psi|^{2}$, а циркуляция зависит от фазы. Например, поле $\psi=F(r)\exp(i q \theta)\exp(-i \Omega t)$ имеет осесимметричную плотность, но по \eqref{eq:29} несёт ток $j_{\theta}=(\hbar q/mr)|F(r)|^{2}$. Поэтому исчезновение лепестков совместимо с ненулевым вращательным потоком.

В расчёте конечное почти круговое поле имеет суммарный заряд \mbox{$q=-3$} и угловой момент около $-3 \hbar$ на частицу. Индекс $m=6$ относится к модуляции плотности и сам по себе не определяет заряд $q$. Его полная временная история не отслеживалась, поэтому установлен именно конечный заряд, а не его сохранение на всей траектории.

Затухание происходит в совместной динамике поля и резервуара с усилением, истощением, потерями и переносом частиц. $A_{6}$ количественно подтверждает исчезновение модуляции, но отдельное модовое разложение не проводилось: приписывать шесть лепестков конкретной паре интерферирующих мод по этим данным нельзя.

\clearpage

\section{Заключение}\label{sec:10}

В однородной модели порог $P_{\rm th}=\gamma_{C} \gamma_{R}/R$ следует из баланса усиления и потерь. Выше него резервуар закрепляется на $n_{R}^{*}=\gamma_{C}/R$, а плотность конденсата равна $(P-P_{\rm th})/\gamma_{C}$. Ненулевая ветвь устойчива к малым однородным плотностным возмущениям; для выбранных параметров колебательная релаксация возникает при $1.0801<s<13.4799$.

Адиабатическое исключение сохраняет стационарные плотности, но может терять запаздывание резервуара. Дополнительное разложение по малой плотности даёт кубическую модель Стюарта--Ландау со стационарной ошибкой $1-1/s$ для полученных коэффициентов. Быстрая подстройка резервуара и малая плотность --- разные условия применимости.

При локализованной накачке порог зависит от пространственной формы моды. Для пятна и кольца он превышает однородный ориентир. Плотность, ток и импульсный спектр показывают совместное действие усиления, резервуарного отталкивания и распространения поля.

При кольцевой накачке шестилепестковая модуляция затухает, оставляя почти осесимметричную плотность с суммарным зарядом \mbox{$q=-3$} и угловым моментом около $-3 \hbar$ на частицу. Конечное состояние воспроизводится при проверенных изменениях дискретизации с той же затравкой. Это конкретное решение модели: статистика выбора заряда, внутренняя структура ядра и устойчивость к произвольным возмущениям не исследованы.

\clearpage

\begin{thebibliography}{9}
\bibitem{ref1}
M.~Wouters and I.~Carusotto.
Excitations in a nonequilibrium Bose--Einstein condensate of exciton polaritons.
\textit{Physical Review Letters} \textbf{99}, 140402 (2007).
\url{https://doi.org/10.1103/PhysRevLett.99.140402}
\bibitem{ref2}
N.~Bobrovska and M.~Matuszewski.
Adiabatic approximation and fluctuations in exciton-polariton condensates.
\textit{Physical Review B} \textbf{92}, 035311 (2015).
\url{https://doi.org/10.1103/PhysRevB.92.035311}
\bibitem{ref3}
S.~Alyatkin et al.
Antiferromagnetic Ising model in a triangular vortex lattice of quantum fluids of light.
\textit{Science Advances} \textbf{10}, eadj1589 (2024).
\url{https://doi.org/10.1126/sciadv.adj1589}
\bibitem{ref4}
S.~N.~Alperin and N.~G.~Berloff.
Multiply charged vortex states of polariton condensates.
\textit{Optica} \textbf{8}, 301--307 (2021).
\url{https://doi.org/10.1364/OPTICA.418377}
\end{thebibliography}

\clearpage

\appendix

\renewcommand{\thesection}{Приложение А.}
\renewcommand{\thesubsection}{А.\arabic{subsection}}

\section{Численная реализация и проверки}\label{sec:A}

\subsection{Начальное поле}\label{sec:A.1}

При $n_{R}(r,0)=0$ поле задаётся одной реализацией суммы плоских волн:

\[
\begin{aligned}\psi(\mathbf r,0)&=\frac{\sqrt{n_{\rm seed}f(r)}}{\sqrt M}\sum_{j=1}^{M}\exp[i(k_{x,j}x+k_{y,j}y+\varphi_j)],\\f(r)&=\frac{P(r)}{P_0P_{\rm th}},\qquad n_{\rm seed}=10^{-6}\;\mu\mathrm m^{-2},\qquad M=24.\end{aligned}
\]

Фазы $\varphi_{j}$ независимы и равномерны на $[0,2\pi)$. Компоненты $k_{x,j}$ и $k_{y,j}$ имеют нормальные распределения с нулевым средним и стандартным отклонением $1\,\mu\mathrm m^{-1}$. Зерно генератора --- $20261008$. Деление на $\sqrt{M}$ сохраняет средний масштаб плотности при изменении числа волн; в отдельной реализации плотность модулируется интерференцией.

При изменении сетки сохраняются волновые векторы и фазы, а та же непрерывная формула вычисляется в новых узлах. Случайность присутствует только в начальном условии. Такая затравка запускает детерминистскую динамику, но не является микроскопической моделью флуктуаций; её параметры не определяются уравнениями резервуара. Чувствительность конечного заряда к выбору реализации здесь не проверялась.

\subsection{Схема интегрирования и спектр}\label{sec:A.2}

Один шаг Стрэнга состоит из кинетического полушага, полного локального шага и второго кинетического полушага. Кинетика вычисляется через FFT с множителем \eqref{eq:25}, локальная система для $\psi$ и $n_{R}$ --- методом Рунге--Кутты четвёртого порядка. Некоммутативность этих частей ограничивает порядок полной схемы вторым.

В расчёте использованы $\hbar=0.6582119569\,\mathrm{meV}\,\mathrm{ps}$ и
\[
K=\hbar^{2}/(2m)=0.676728618\,\mathrm{meV}\,\mu\mathrm m^{2}.
\] Округлённые значения в основном тексте служат для представления результатов.

Для конвенции \eqref{eq:32} дискретное преобразование имеет вид $\widetilde\psi\simeq \Delta x \Delta y$ FFT$2(\psi)$. Множитель площади ячейки согласует сумму с интегралом и сохраняет физический масштаб при изменении сетки. Размерности $\widetilde\psi$ и $I_{k}$ --- $\mu\mathrm m$ и $\mu\mathrm m^{2}$ соответственно. В логарифме на рис.~\ref{fig:7} интенсивность делится на $1\,\mu\mathrm m^{2}$.

\subsection{Критерии установления режима}\label{sec:A.3}

Анализируются полные окна длительностью $100\,\mathrm{ps}$. Населённости записываются через $1\,\mathrm{ps}$, пространственные поля --- через $5\,\mathrm{ps}$. Поля сравниваются при задержках $20$ и $35\,\mathrm{ps}$; изменения относительно начала окна дополнительно накапливаются с шагом $1\,\mathrm{ps}$.

\begin{center}\begin{tabular}{>{\raggedright\arraybackslash}p{0.72\textwidth}>{\raggedright\arraybackslash}p{0.17\textwidth}}\toprule

\textbf{Контролируемая величина} & \textbf{Допуск}\\\midrule

Относительный размах каждой населённости & $0.5\%$\\

Изменения плотностей конденсата и резервуара & $1\%$\\

Изменение поля после согласования общей фазы & $1\%$\\

Изменение тока, где критерий применим & $2\%$\\

\bottomrule\end{tabular}\end{center}

Согласование общей фазы устраняет различие, вызванное только множителем $\exp(-i \Omega t)$. Для кольца отдельно проверяется форма плотности в $\mathrm{core}$: $|x|,|y|\le 40\,\mu\mathrm m$. Определяются

\[
\begin{aligned}u(\mathbf r,t)&=\frac{n_C(\mathbf r,t)}{\displaystyle\int_{\rm core}n_C(\mathbf r,t)\,\mathrm d^2r},\\D_u(t)&=\frac{\|u(t)-u(t_0)\|_2}{\|u(t_0)\|_2},\\\|v\|_2^2&=\int_{\rm core}|v|^2\,\mathrm d^2r.\end{aligned}
\]

Приводится максимальное $D_{u}$ за окно. Нормировка отделяет изменение формы от изменения числа частиц внутри $\mathrm{core}$. Для окна $1900\text{--}2000\,\mathrm{ps} \max D_{u}=0.0234\%$, а $(\max N_{C}-\min N_{C})/\operatorname{mean} N_{C}=1.65 \times 10^{-6}\%$. Проверки поздних колебаний не выявили значимого самоподдерживающегося периода. Эти величины описывают изменение одного решения во времени, а не его ошибку относительно независимого расчёта.

\subsection{Дискретизация, границы и баланс}\label{sec:A.4}

Для гауссова пятна проверены сетки $256^{2}, 512^{2}$ и $1024^{2}$, шаги $0.05, 0.025$ и $0.0125\,\mathrm{ps}$ и два размера области. Наблюдаемый временной порядок равен $1.998$. Основные расхождения населённостей находятся в пределах $0.5\%$. Одна проверка границ для меньшей области имеет предупреждение: сходство спектральных карт само по себе его не отменяет.

Для кольца эталон --- $512^{2}, \Delta t=0.05\,\mathrm{ps}, L=256\,\mu\mathrm m, t_{\rm final}=2000\,\mathrm{ps}$. Дополнительные запуски выполнены независимо от $t=0$ с тем же начальным полем. Ниже все расхождения приведены в процентах относительно эталона.

\begin{center}\begin{tabular}{>{\raggedright\arraybackslash}p{0.32\textwidth}>{\raggedright\arraybackslash}p{0.16\textwidth}>{\raggedright\arraybackslash}p{0.16\textwidth}>{\raggedright\arraybackslash}p{0.22\textwidth}}\toprule

\textbf{Проверка} & \textbf{$N_{C}$} & \textbf{$N_{R}$} & \textbf{Плотность $n_{C}$, относительная норма $L^{2}$}\\\midrule

Сетка $256^{2}, \Delta t=0.05\,\mathrm{ps}$ & Не более $0.000031\%$ & Не более $0.000031\%$ & Не более $0.000031\%$\\

Сетка $512^{2}, \Delta t=0.025\,\mathrm{ps}$ & $0.00594\%$ & $0.000950\%$ & $0.00559\%$\\

\bottomrule\end{tabular}\end{center}

В пространственной строке указан общий верхний предел трёх расхождений, а не три одинаковых измеренных значения. Плотности сравниваются по всей области; это другая метрика, чем нормированная форма $u$ в \ref{sec:A.3}. Сетка $256^{2}$ грубее эталонной, поэтому сравнение подтверждает согласованность двух разрешений, но не устанавливает пространственный порядок сходимости.

Оба расчёта достигают установившегося режима и воспроизводят \mbox{$q=-3$} на пяти общих допустимых контурах, а также $L_{z}/(\hbar N_{C})\simeq -3$. Конечное кольцевое состояние проходит проверки баланса \eqref{eq:27}, плотности у краёв и высокочастотных компонент спектра. Последняя проверка контролирует вклад компонент около предела пространственного разрешения.

\subsection{Диагностика циркуляции и угловой модуляции}\label{sec:A.5}

Для конечного кольцевого поля используются контуры радиусами $4, 6, 7.5, 9$ и $12\,\mu\mathrm m$. Обход направлен против часовой стрелки; на каждом контуре получается \mbox{$q=-3$}. Фаза в областях малой плотности на рисунке скрыта.

В \eqref{eq:38} множитель $2$ выбран так, что для $n_{C}=n_{0}(r)[1+a \cos(m \theta+\alpha)]$ при полном угловом интегрировании $A_{m}=|a|$. В расчёте интеграл ограничен квадратом $\mathrm{core}$, поэтому $A_{m}$ трактуется как нормированная проекция плотности на угловую гармонику. Значения $29.1\%, 0.567\%$ и $0.0107\%$ относятся к $m=6$ при $1000, 1500$ и $2000\,\mathrm{ps}$.

\subsection{Выделение ведущей моды при поиске порога}\label{sec:A.6}

Линейное уравнение \eqref{eq:36} интегрируется при фиксированном $n_{R,0}=P/\gamma_{R}$. Скорость роста оценивается по двум поздним окнам по $100\,\mathrm{ps}$ с проверкой качества оценки и перекрытий пространственных мод. Если ведущая мода выделяется медленно, расчёт продлевается: смесь мод на раннем участке может давать другой наклон $\ln N_{\rm lin}$.

Чтобы избежать переполнения или потери малых чисел, поле периодически перенормируется, а логарифмы множителей накапливаются. Линейность позволяет восстановить скорость роста без изменения формы моды. Используются повторные случайные затравки и изменения сетки и размера области. Это проверки линейного порога, а не статистики конечного нелинейного заряда.

При сопоставлении с разделом $3$ учитывается, что \eqref{eq:37} описывает рост амплитуды, а собственные значения якобиана --- рост плотностей. Кроме того, максимум на рис.~\ref{fig:1} далеко ниже порога может соответствовать резервуарной моде; всю эту кривую нельзя отождествить с $\sigma$ простым делением на два.

\clearpage

% Приложение Б: русская версия приложения B английского отчёта.
\renewcommand{\thesection}{Приложение Б.}
\renewcommand{\thesubsection}{Б.\arabic{subsection}}
\renewcommand{\theequation}{Б.\arabic{equation}}
\renewcommand{\theHequation}{B.\arabic{equation}}
\setcounter{equation}{0}
\section{Уравнения в пространстве Фурье и схема расщепления}\label{sec:B}

\subsection{Что меняется при преобразовании Фурье}\label{sec:B.1}
Для конвенции преобразования Фурье \eqref{eq:32} пространственное дифференцирование даёт
\(\mathcal F[\nabla_\perp^2\psi]=-k^2\widetilde\psi\), где
\(k^2=k_x^2+k_y^2\). После преобразования полная система принимает вид
\begin{align}
\partial_t\widetilde\psi={}&-\frac{iKk^2}{\hbar}\widetilde\psi
-\frac{\gamma_C}{2}\widetilde\psi
+\frac{R}{2}\mathcal F[n_R\psi]\notag\\
&-\frac{i}{\hbar}\left(g_C\mathcal F[|\psi|^2\psi]
+g_R\mathcal F[(n_R+\eta P)\psi]\right),\label{eq:B.field}\\
\partial_t\widetilde n_R={}&\widetilde P-\gamma_R\widetilde n_R
-R\mathcal F[n_R|\psi|^2].\label{eq:B.reservoir}
\end{align}
Только кинетическое слагаемое диагонально по волновому вектору. Произведения в координатном пространстве переходят в свёртки в пространстве Фурье, поэтому полная динамика не распадается на независимые уравнения для каждого волнового вектора. В расчёте кинетическая эволюция выполняется в пространстве Фурье, а локальные слагаемые вычисляются в координатном пространстве.

\subsection{Кинетические полушаги}\label{sec:B.2}
На кинетическом шаге поле подчиняется уравнению
\begin{equation}
\begin{aligned}
\partial_t\widetilde\psi&=-\frac{iKk^2}{\hbar}\widetilde\psi,\\
\widetilde\psi(t+\tau)&=e^{-iKk^2\tau/\hbar}\widetilde\psi(t),
\qquad \partial_t n_R=0.
\end{aligned}\label{eq:B.kinetic}
\end{equation}
Соответствующий точный дискретный оператор кинетической эволюции имеет вид
\begin{equation}
\mathcal T_\tau\psi=
\operatorname{IFFT2}\!\left[
e^{-iK(k_x^2+k_y^2)\tau/\hbar}
\operatorname{FFT2}(\psi)\right].\label{eq:B.kinetic-map}
\end{equation}
Обозначим через \(\mathcal L_\tau\) совместную локальную эволюцию поля и резервуара, определённую ниже. Один шаг реализованной схемы Стрэнга записывается как
\begin{equation}
\begin{aligned}
(\psi^{(1)},n_R^{(1)})&=(\mathcal T_{\Delta t/2}\psi^n,n_R^n),\\
(\psi^{(2)},n_R^{(2)})&=\mathcal L_{\Delta t}
(\psi^{(1)},n_R^{(1)}),\\
(\psi^{n+1},n_R^{n+1})&=(\mathcal T_{\Delta t/2}\psi^{(2)},n_R^{(2)}).
\end{aligned}\label{eq:B.strang}
\end{equation}
На каждом кинетическом полушаге выполняются прямое FFT, умножение на фазовый множитель и обратное FFT. Порядок операций: кинетический полушаг, полный локальный шаг и второй кинетический полушаг. Локальные коэффициенты изменяются вместе с полем и резервуаром, поэтому локальный шаг учитывает их совместную динамику.

\subsection{Локальный шаг}\label{sec:B.3}
В каждом узле координатной сетки локальный шаг решает систему
\begin{equation}
\begin{aligned}
\dot\psi&=\left[\frac{Rn_R-\gamma_C}{2}
-\frac{i}{\hbar}\bigl(g_C|\psi|^2+g_R(n_R+\eta P)\bigr)\right]\psi,\\
\dot n_R&=P-(\gamma_R+R|\psi|^2)n_R.
\end{aligned}\label{eq:B.local}
\end{equation}
Запишем систему как \(\dot Y=f(Y;P)\), где
\(Y=(\operatorname{Re}\psi,\operatorname{Im}\psi,n_R)^T\).
В нелинейных расчётах конденсата и резервуара локальный шаг вычисляется методом Рунге--Кутты четвёртого порядка:
\begin{equation}
\begin{aligned}
a_1&=f(Y;P),& a_2&=f(Y+\tfrac{\Delta t}{2}a_1;P),\\
a_3&=f(Y+\tfrac{\Delta t}{2}a_2;P),& a_4&=f(Y+\Delta t\,a_3;P),\\
Y_{\mathrm{out}}&=Y+\frac{\Delta t}{6}(a_1+2a_2+2a_3+a_4).
\end{aligned}\label{eq:B.rk4}
\end{equation}
Плотность \(n_C=|\psi|^2\) и коэффициенты, зависящие от резервуара, пересчитываются на каждой стадии; \(n_C\) не интегрируется как отдельная переменная. Узлы эволюционируют независимо только внутри локального шага. Кинетические шаги связывают поле в разных точках пространства. Хотя кинетическая эволюция точна для дискретных фурье-мод, а локальный интегратор имеет четвёртый порядок, полная схема расщепления имеет второй порядок по времени.


При линейном поиске порога резервуар фиксирован на значении
\(n_{R,0}=P/\gamma_R\), а кубическое слагаемое поля отброшено, как в \eqref{eq:36}. Локальный коэффициент тогда не зависит от времени, и локальный шаг вычисляется точной экспонентой:
\begin{equation}\label{eq:B.linear-local}
\psi_{\mathrm{out}}=\exp\!\left\{\Delta t\left[
\frac{RP/\gamma_R-\gamma_C}{2}
-\frac{i g_R}{\hbar}\left(\frac{P}{\gamma_R}+\eta P\right)
\right]\right\}\psi_{\mathrm{in}}.
\end{equation}
До и после этого локального шага выполняются те же кинетические полушаги. Экспонента точна для локальной задачи с фиксированным резервуаром; расщепление пространственных операторов по-прежнему ограничивает порядок полной схемы вторым.

\subsection{Пространственная сетка и проверки точности}\label{sec:B.4}
Поле \(\psi\) и плотность резервуара \(n_R\) заданы на одной равномерной периодической декартовой сетке. Накачка --- фиксированный массив на этой сетке, а \(n_C\) вычисляется в каждом узле. Для используемых чётных размеров сетки координаты и волновые векторы FFT имеют вид
\begin{equation}
\begin{aligned}
x_j&=-\frac{L_x}{2}+j\Delta x,
&\Delta x&=\frac{L_x}{N_x},\quad j=0,\ldots,N_x-1,\\
y_l&=-\frac{L_y}{2}+l\Delta y,
&\Delta y&=\frac{L_y}{N_y},\quad l=0,\ldots,N_y-1,\\
k_{x,m}&=\frac{2\pi m}{L_x},
&m&=0,\ldots,\frac{N_x}{2}-1,-\frac{N_x}{2},\ldots,-1,
\end{aligned}\label{eq:B.grid}
\end{equation}
Для \(k_y\) используется аналогичный порядок. Верхняя граница координатного интервала не включается повторно. Нормировочные множители прямого и обратного FFT сокращаются в \eqref{eq:B.kinetic-map}.

\begin{center}
\small\begin{tabular}{lrrrr}
\hline
Расчёт & Сетка & \(L\) (\(\mu\mathrm m\)) & \(\Delta x\) (\(\mu\mathrm m\)) & \(\Delta t\) (\(\mathrm{ps}\))\\
\hline
Гауссово пятно & \(512^2\) & 256 & 0.5 & 0.025\\
Кольцо & \(512^2\) & 256 & 0.5 & 0.05\\
Линейный поиск порога & \(128^2\) & 128 & 1 & 0.05\\
\hline
\end{tabular}
\end{center}
В основных нелинейных расчётах масштаб накачки \(d=2.5\,\mu\mathrm m\) соответствует пяти пространственным шагам, а расчётная область значительно больше области накачки. Эти соотношения объясняют выбор сетки, но сами по себе не устанавливают точность. Проверки сетки, временного шага, границ и спектра приведены в приложении~\ref{sec:A.4}. Границы линейных пороговых интервалов дополнительно проверены на сетке \(256^2\) при \(L=128\,\mu\mathrm m\) и на сетке \(192^2\) при \(L=192\,\mu\mathrm m\). Пространственное разрешение и размер области проверяются отдельно.

\end{document}
